
Model accuracy does not capture the complete picture of model behavior. Analysis of the Small CNN Zoo---a collection of 193 convolutional neural networks trained on CIFAR-10 with varying hyperparameters and random seeds---reveals that 67.6\% of class-wise accuracy variance remains unexplained by overall model accuracy. This finding exceeds the detection threshold of 0.05 by a factor of 13.5, establishing that trained models develop distinct behavioral fingerprints: class-specific performance patterns that cannot be reduced to a single accuracy number.

This observation has practical implications. Model selection typically relies on aggregate metrics, yet two models with identical 85\% accuracy may fail on entirely different classes. Understanding these behavioral differences from weights alone---without requiring inference---could enable efficient model auditing, ensemble construction, and failure mode prediction. The central question is whether this behavioral structure can be extracted from weight matrices.

\subsubsection{The Gap in Weight-Space Learning}

Recent advances in weight-space learning have achieved success at scalar property prediction. SANE achieves $R^2$ > 0.9 for predicting overall accuracy from weight representations on small CNNs. Unterthiner et al. demonstrate that simple weight statistics (means, standard deviations, spectral norms) can predict accuracy with $R^2$ approaching 0.97. Neural Functionals provide permutation-equivariant architectures for processing weights while respecting hidden neuron symmetries.

However, these methods predict scalar properties. The structured prediction problem---forecasting a 10-dimensional class-wise accuracy vector rather than a single number---remains unexplored. More fundamentally, the question of whether behavioral information (beyond accuracy) is encoded in weights at all has not been systematically tested.

\subsubsection{Contributions}

This work addresses two foundational questions:

\begin{enumerate}
\item \textbf{Existence (H-E1):} Does meaningful behavioral variance exist in model zoos beyond what overall accuracy explains?
\item \textbf{Mechanism (H-M1):} Can simple weight-space features extract this behavioral signal?
\end{enumerate}

For existence, decomposition of class-wise accuracy variance into components explained by overall accuracy, per-class baseline difficulty, and residual behavioral variance yields a residual ratio of 0.676---67.6\% of variance is not explained by the stratified baseline---establishing that behavioral fingerprints exist in model zoos.

For mechanism, per-layer weight statistics (mean, standard deviation, minimum, maximum, L2 norm across 5 layers = 25 features) are tested for predicting class-wise accuracy profiles. Using Ridge regression ($\alpha$ = 1.0) with 80/20 train/test splits (154 training, 39 test models), weight features achieve mean $R^2$ = -0.08 compared to the stratified baseline mean $R^2$ = 0.12. The mechanism test fails: simple weight statistics do not capture behavioral variance better than knowing overall accuracy alone.

This negative result is informative. It establishes that while behavioral fingerprints exist, extracting them from weights using simple statistics is non-trivial.
